% problem_id: S001-lemma-brown-bis
% source_id: S001 (Stacks Project)
% source_file: derived.tex
% statement_locator: derived.tex lines 11468--11486
% proof_locator: derived.tex lines 11488--11631
% env: lemma
% id_note: label 在本来源内唯一
% pairing: tight (verified)
% status: complete (statement + author proof verbatim + reason + locators)
%
\begin{lemma}
\label{lemma-brown-bis}
\begin{reference}
Weak version of \cite[Theorem A]{Krause}
\end{reference}
Let $\mathcal{D}$ be a triangulated category with direct sums.
Suppose given a set $\mathcal{E}$ of objects of $\mathcal{D}$ such that
\begin{enumerate}
\item if $X$ is a nonzero object of $\mathcal{D}$, then there exists
an $E \in \mathcal{E}$ and a nonzero map $E \to X$, and
\item given objects $X_n$, $n \in \mathbf{N}$ of $\mathcal{D}$,
$E \in \mathcal{E}$, and $\alpha : E \to \bigoplus X_n$,
there exist $E_n \in \mathcal{E}$ and $\beta_n : E_n \to X_n$ and a morphism
$\gamma : E \to \bigoplus E_n$ such that
$\alpha = (\bigoplus \beta_n) \circ \gamma$.
\end{enumerate}
Let $H : \mathcal{D} \to \textit{Ab}$ be a contravariant cohomological functor
which transforms direct sums into products. Then $H$ is representable.
\end{lemma}

\begin{proof}
This proof is very similar to the proof of Lemma \ref{lemma-brown}.
We may replace $\mathcal{E}$ by
$\bigcup_{i \in \mathbf{Z}} \mathcal{E}[i]$ and assume
that $\mathcal{E}$ is preserved by shifts.
Consider pairs $(E, a)$ where $E \in \mathcal{E}$ and $a \in H(E)$ and set
$$
X_1 = \bigoplus\nolimits_{(E, a)} E
$$
Since $H(X_1) = \prod_{(E, a)} H(E)$ we see that $(a)_{(E, a)}$
defines an element $a_1 \in H(X_1)$. Set $H_1 = \Hom_\mathcal{D}(- , X_1)$.
By Yoneda's lemma (Categories, Lemma \ref{categories-lemma-yoneda})
the element $a_1$ defines a natural transformation $H_1 \to H$.

\medskip\noindent
We are going to inductively construct $X_n$ and transformations
$a_n : H_n \to H$ where $H_n = \Hom_\mathcal{D}(-, X_n)$.
Namely, we apply the procedure
above to the functor $\Ker(H_n \to H)$ to get an object
$$
K_{n + 1} = \bigoplus\nolimits_{(E, k),\ k \in \Ker(H_n(E) \to H(E))} E
$$
and a transformation $\Hom_\mathcal{D}(-, K_{n + 1}) \to \Ker(H_n \to H)$.
By Yoneda's lemma the composition $\Hom_\mathcal{D}(-, K_{n + 1}) \to H_n$
gives a morphism $K_{n + 1} \to X_n$. We choose
a distinguished triangle
$$
K_{n + 1} \to X_n \to X_{n + 1} \to K_{n + 1}[1]
$$
in $\mathcal{D}$. The element $a_n \in H(X_n)$ maps to zero
in $H(K_{n + 1})$ by construction. Since $H$ is cohomological
we can lift it to an element $a_{n + 1} \in H(X_{n + 1})$.

\medskip\noindent
Set $X = \text{hocolim} X_n$. Applying $H$
to the defining distinguished triangle
$$
\bigoplus X_n \to
\bigoplus X_n \to X \to
\bigoplus X_n[1]
$$
we obtain an exact sequence
$$
\prod H(X_n) \leftarrow
\prod H(X_n) \leftarrow
H(X)
$$
Thus there exists an element $a \in H(X)$ mapping to $(a_n)$
in $\prod H(X_n)$. Hence a natural transformation
$\Hom_\mathcal{D}(- , X) \to H$ such that
$$
\Hom_\mathcal{D}(-, X_1) \to
\Hom_\mathcal{D}(-, X_2) \to
\Hom_\mathcal{D}(-, X_3) \to \ldots \to
\Hom_\mathcal{D}(-, X) \to H
$$
commutes. We claim that $\Hom_\mathcal{D}(-, X) \to H(-)$ is an isomorphism.

\medskip\noindent
Let $E \in \mathcal{E}$. Let us show that
$$
\Hom_\mathcal{D}(E, \bigoplus X_n) \to  \Hom_\mathcal{D}(E, \bigoplus X_n)
$$
is injective. Namely, let $\alpha : E \to \bigoplus X_n$. Then
by assumption (2) we obtain a factorization
$\alpha = (\bigoplus \beta_n) \circ \gamma$.
Since $E_n \to X_n \to X_{n + 1}$ is zero by construction, we see that
the composition $\bigoplus E_n \to \bigoplus X_n \to \bigoplus X_n$
is equal to $\bigoplus \beta_n$. Hence also the composition
$E \to \bigoplus X_n \to \bigoplus X_n$ is equal to $\alpha$.
This proves the stated injectivity and hence also
$$
\Hom_\mathcal{D}(E, \bigoplus X_n[1]) \to \Hom_\mathcal{D}(E, \bigoplus X_n[1])
$$
is injective. It follows that we have an exact sequence
$$
\Hom_\mathcal{D}(E, \bigoplus X_n) \to
\Hom_\mathcal{D}(E, \bigoplus X_n) \to
\Hom_\mathcal{D}(E, X) \to 0
$$
for all $E \in \mathcal{E}$.

\medskip\noindent
Let $E \in \mathcal{E}$ and let $f : E \to X$ be a morphism.
By the previous paragraph, we may choose $\alpha : E \to \bigoplus X_n$
lifting $f$. Then by assumption (2) we obtain a factorization
$\alpha = (\bigoplus \beta_n) \circ \gamma$.
For each $n$ there is a morphism $\delta_n : E_n \to X_1$ such that
$\delta_n$ and $\beta_n$ map to the same element of $H(E_n)$. Then
the compositions
$$
E_n \to X_n \to X_{n + 1}
\quad\text{and}\quad
E_n \to X_1 \to X_{n + 1}
$$
are equal by construction of $X_n \to X_{n + 1}$. It follows that
$$
\bigoplus E_n \to \bigoplus X_n \to X
\quad\text{and}\quad
\bigoplus E_n \to \bigoplus X_1 \to X
$$
are the same too. Observing that $\bigoplus X_1 \to X$ factors
as $\bigoplus X_1 \to X_1 \to X$, we conclude that
$$
\Hom_\mathcal{D}(E, X_1) \to \Hom_\mathcal{D}(E, X)
$$
is surjective. Since by construction the map
$\Hom_\mathcal{D}(E, X_1) \to H(E)$ is surjective and by
construction the kernel of this map is annihilated by
$\Hom_\mathcal{D}(E, X_1) \to \Hom_\mathcal{D}(E, X)$
we conclude that $\Hom_\mathcal{D}(E, X) \to H(E)$ is
a bijection for all $E \in \mathcal{E}$.

\medskip\noindent
To finish the proof, consider the subcategory
$$
\mathcal{D}' =
\{Y \in \Ob(\mathcal{D}) \mid \Hom_\mathcal{D}(Y[n], X) \to H(Y[n])
\text{ is an isomorphism for all }n\}
$$
As $\Hom_\mathcal{D}(-, X) \to H$ is a transformation between
cohomological functors,
the subcategory $\mathcal{D}'$ is a strictly full, saturated, triangulated
subcategory of $\mathcal{D}$ (details omitted; see proof of
Lemma \ref{lemma-homological-functor-kernel}). Moreover, as both
$H$ and $\Hom_\mathcal{D}(-, X)$ transform direct sums into products,
we see that direct sums of objects of $\mathcal{D}'$ are in $\mathcal{D}'$.
Thus derived colimits of objects of $\mathcal{D}'$ are in $\mathcal{D}'$.
Since $\mathcal{E}$ is preserved by shifts, we conclude that
$\mathcal{E} \subset \Ob(\mathcal{D}')$ by the result of the
previous paragraph. To finish the proof we have to show that
$\mathcal{D}' = \mathcal{D}$.

\medskip\noindent
Let $Y$ be an object of $\mathcal{D}$ and set $H(-) = \Hom_\mathcal{D}(-, Y)$.
Then $H$ is a cohomological functor which transforms direct sums into products.
By the construction in the first part of the proof we obtain a morphism
$\colim X_n = X \to Y$ such that
$\Hom_\mathcal{D}(E, X) \to \Hom_\mathcal{D}(E, Y)$ is bijective for all
$E \in \mathcal{E}$. Then assumption (1) tells us that $X \to Y$ is an
isomorphism! On the other hand, by construction $X_1, X_2, \ldots$
are in $\mathcal{D}'$ and so is $X$. Thus $Y \in \mathcal{D}'$ and
the proof is complete.
\end{proof}
